%%==========================================================================
\section{Computational Study}
\label{sec:computational}
%%==========================================================================

\subsection{Implementation and environment}

All models are implemented in Python using Pyomo. The reported exact solves
use unrestricted IBM ILOG CPLEX~12.6 for AMPL through Pyomo's
\texttt{cplexamp} interface. The deterministic instance generator and all
benchmark scripts are included in the companion repository. Every reported
solution is checked independently by replaying the aircraft route and its
maintenance counter.

\subsection{Cumulative-flight-time row verification}

The original cumulative-flight-time row and the natural endpoint-split
replacement were evaluated on the hand-built four-day counterexample
\texttt{data/instances/c13\_loophole\_test.json}. The instance fixes one
aircraft, $T^A_{\max}=600$ minutes, $M_{14}=1200$ minutes, 1,200 minutes of
flying in $(1,4]$, and two check patterns: checks at both endpoints, and a
check only at the end. The actual production row objects admit complementary
violating patterns: the original row is relaxed when only one endpoint is
checked, while the split replacement is relaxed when both endpoints carry
checks. The exact cumulative-state and ordered event-state formulations reject
both patterns. This directly confirms the algebraic analysis in
\Cref{subsec:c11_flaw,subsec:c11_corrected}.

\begin{table}[t]
\centering\small
\caption{Four-way numeric verification of the cumulative-flight-time flaw.
The fixed counterexample has $T^A_{\max}=600$, $M_{14}=1200$, and LHS
$\sum_{i\in F_{1,4}}x_{ij}\tilde t_i=1200$. ``Accepts'' means every row of
that formulation is satisfied; ``rejects'' means the violation is detected.}
\label{tab:c13_loophole}
\begin{tabular}{lcccc}
\toprule
Scenario & Khaled~(13) & Legacy split & Legacy corrected strengthened & Event-based optimized \\
\midrule
Both endpoints checked & rejects & accepts (wrong) & rejects & rejects \\
Only end checked & accepts (wrong) & rejects & rejects & rejects \\
\bottomrule
\end{tabular}
\end{table}

\subsection{Four-method type-A comparison}
\label{subsec:event_measurements}

The ordered implementation removes explicit successor variables and propagates
an aircraft-local prefix state through flights sorted by departure time. A
maintenance event is attached to its triggering flight, and the pre-check row
prevents an event from repairing an already overdue arrival. The four methods
evaluated in the matched A-only experiment are the original Khaled row,
the endpoint split, the strengthened corrected legacy rows, and the optimized
ordered event-based formulation. Their builder identifiers are respectively
\texttt{legacy\_paper\_c13}, \texttt{legacy\_endpoint\_split},
\texttt{legacy\_corrected\_strengthened}, and \texttt{event\_based\_optimized}. A result is described
as objective-matched only when both
compared methods are optimal.

\subsubsection{A-only timing versus planning horizon}

We compared the four methods on the same corrected feasible instances with
$|P|=10$ and $h\in\{7,15,21,30\}$ using unrestricted CPLEX~12.6 from
\texttt{F:/Progs/IBM.ILOG.CPLEX.for.AMPL.v12.6-EAT/CPLEXamp.exe}.
The instances use overnight returns and initial A/B hours chosen so that the
seeded rotations are feasible under day/night maintenance semantics.
The complete measurements are in
\texttt{results/tables/formulation\_a\_comparison\_four\_strengthened.csv}.
Thus this table is a one-instance-per-horizon experiment, not a reproduction
of Khaled et al.'s Table~6. The four input files contain 36 flights for $h=7$
and 40 flights for each of $h\in\{15,21,30\}$, with $|P|=10$, an A-check
threshold of 270 minutes, initial A-hours of 90 minutes, 90-minute flight
legs, and maintenance capacity 10 at MRO. They were generated by
\texttt{src/generate\_feasible\_instances.py} with
\texttt{--maintenance-families A}; the same JSON file is passed to all four
builders for each horizon.

\begin{table}[t]
\centering\scriptsize
\caption{Four-method A-only comparison by planning horizon on the corrected
feasible family. All reported runs are optimal. The method names identify the
exact builder and C13 treatment.}
\label{tab:a_only_horizon_timing}
\resizebox{\linewidth}{!}{%
\begin{tabular}{r l l r r r r r}
\toprule
$h$ & Method & Status & Objective & CPU (s) & Wall (s) & Vars & Cons. \\
\midrule
7  & Legacy paper & optimal & 17526 & 0.326 & 0.673 & 1720 & 3035 \\
  & Legacy split & optimal & 17526 & 0.341 & 0.714 & 1720 & 3315 \\
  & Legacy corrected strengthened & optimal & 17526 & 0.334 & 0.671 & 1810 & 2255 \\
  & Event-based optimized & optimal & 17526 & 0.317 & 0.665 & 900 & 2816 \\
15 & Legacy paper & optimal & 19586 & 0.394 & 0.815 & 2520 & 4304 \\
  & Legacy split & optimal & 19586 & 0.433 & 0.893 & 2520 & 5084 \\
  & Legacy corrected strengthened & optimal & 19586 & 0.354 & 0.716 & 2660 & 3424 \\
  & Event-based optimized & optimal & 19586 & 0.330 & 0.684 & 1000 & 3150 \\
21 & Legacy paper & optimal & 19570 & 0.481 & 0.895 & 2520 & 4304 \\
  & Legacy split & optimal & 19570 & 0.438 & 0.892 & 2520 & 5084 \\
  & Legacy corrected strengthened & optimal & 19570 & 0.372 & 0.736 & 2660 & 3424 \\
  & Event-based optimized & optimal & 19570 & 0.337 & 0.689 & 1000 & 3150 \\
30 & Legacy paper & optimal & 19592 & 0.389 & 0.813 & 2520 & 4304 \\
  & Legacy split & optimal & 19592 & 0.418 & 0.878 & 2520 & 5084 \\
  & Legacy corrected strengthened & optimal & 19592 & 0.358 & 0.737 & 2660 & 3424 \\
  & Event-based optimized & optimal & 19592 & 0.323 & 0.696 & 1000 & 3150 \\
\bottomrule
\end{tabular}
}%
\end{table}

All four methods are optimal on every horizon in the corrected feasible family.
All four methods have identical objectives here, confirming objective parity
after the event-reset correction. The optimized event-based method uses fewer variables
than every legacy variant. It is faster than the paper and split variants on
all four horizons and faster than corrected legacy on the three longer
horizons; strengthened corrected legacy is slightly faster at $H=7$. These results support
a performance comparison under common feasibility semantics without implying
universal runtime dominance.

\begin{figure}[t]
\centering
\includegraphics[width=0.86\linewidth]{figures/ordered_a_wall_vs_horizon_extended_p20.png}
\caption{A-only solve wall time versus planning horizon for the four-method
comparison at $|P|=20$ in the extended suite.}
\label{fig:ordered_a_wall_horizon}
\end{figure}

\begin{figure}[t]
\centering
\includegraphics[width=0.86\linewidth]{figures/ordered_a_cpu_vs_horizon_extended_p20.png}
\caption{A-only solve CPU time versus planning horizon for the four-method
comparison at $|P|=20$ in the extended suite.}
\label{fig:ordered_a_cpu_horizon}
\end{figure}

\subsubsection{A-only CPU scaling versus fleet size}

The fleet-size experiment below supplies the matched four-method comparison
that is needed to complement the horizon study.

\subsubsection{Heavy A-only fleet-size scaling}

To assess how model size changes with the number of aircraft, we use the
A-only fleet-size suite in
\texttt{results/tables/formulation\_a\_fleet\_four\_strengthened.csv}.
These cases use $h=7$, $|P|\in\{5,10,15,20\}$, A checks only, overnight
returns, and one common feasible instance at each fleet size. All four
formulations are run on the same JSON instance for each value of $|P|$:
the Khaled paper row, legacy endpoint split, strengthened corrected legacy, and
optimized ordered Event-based model. All 16 runs are optimal and objective-matched.
The plots report active scalar variables and constraints counted directly from
the built Pyomo models; the constraint plot is the requested dependency-size
view.

Table~\ref{tab:a_fleet_timing} gives the exact CPU and wall-time values used
in Figures~\ref{fig:a_cpu_fleet} and~\ref{fig:a_wall_fleet}. Values are in
seconds; each group of four columns follows the method order in the first
column.

\begin{table}[t]
\centering\scriptsize
\caption{Heavy A-only fleet-size timing data used in Figures 6 and 7.
Values are copied from \texttt{formulation\_a\_fleet\_four\_strengthened.csv}.}
\label{tab:a_fleet_timing}
\resizebox{\linewidth}{!}{%
\begin{tabular}{r r r r r r r r r}
\toprule
$|P|$ & \multicolumn{4}{c}{CPU time (s)} & \multicolumn{4}{c}{Wall time (s)} \\
\cmidrule(lr){2-5}\cmidrule(lr){6-9}
& Paper & Split & Corrected strengthened & Event optimized & Paper & Split & Corrected strengthened & Event optimized \\
\midrule
5  & 0.278806 & 0.284396 & 0.306999 & 0.288996 & 0.563580 & 0.570828 & 0.593329 & 0.570974 \\
10 & 0.322849 & 0.323723 & 0.305997 & 0.319999 & 0.688894 & 0.672230 & 0.633328 & 0.661015 \\
15 & 0.366002 & 0.374803 & 0.365999 & 0.359999 & 0.789226 & 0.827724 & 0.791674 & 0.779326 \\
20 & 0.435350 & 0.458497 & 0.418998 & 0.448000 & 1.016204 & 1.062455 & 0.931895 & 0.999949 \\
\bottomrule
\end{tabular}
}%
\end{table}

\begin{figure}[t]
\centering
\includegraphics[width=0.86\linewidth]{figures/a_variables_vs_fleet_extended_h30.png}
\caption{Active model variables versus fleet size at $H=30$ in the extended
four-method A-only suite.}
\label{fig:a_variables_fleet}
\end{figure}

\begin{figure}[t]
\centering
\includegraphics[width=0.86\linewidth]{figures/a_constraints_vs_fleet_extended_h30.png}
\caption{Active model constraints versus fleet size at $H=30$ in the extended
four-method A-only suite.}
\label{fig:a_constraints_fleet}
\end{figure}

\begin{figure}[t]
\centering
\includegraphics[width=0.86\linewidth]{figures/a_cpu_vs_fleet_extended_h30.png}
\caption{Solver CPU time versus fleet size at $H=30$ in the extended
four-method A-only suite.}
\label{fig:a_cpu_fleet}
\end{figure}

\begin{figure}[t]
\centering
\includegraphics[width=0.86\linewidth]{figures/a_wall_vs_fleet_extended_h30.png}
\caption{Solver wall time versus fleet size at $H=30$ in the extended
four-method A-only suite.}
\label{fig:a_wall_fleet}
\end{figure}

\subsection{Summary and limitations}

The corrected type-A model rejects the cumulative-flight-time counterexample.
The four-way horizon experiment shows that all formulations are feasible and
optimal on the corrected generated family, with the strengthened legacy and
optimized event variants improving on the legacy split in solve time. The
optimized event formulation also uses fewer variables and matches the legacy objectives on this
family. The study does not claim a solver-independent runtime advantage.

\subsubsection{Expanded four-method suite}

To broaden the single-cell horizon and fleet-size snapshots above, we ran the
dedicated four-method suite in
\texttt{results/tables/four\_method\_suite\_strengthened.csv}. It contains 12 generated
A-only performance cases on the Cartesian grid
$|P|\in\{5,10,20,40\}$ and $H\in\{7,15,30\}$, plus the two regression cases
\texttt{c13\_loophole\_test} and
\texttt{event\_vs\_legacy\_clean\_test}. Every performance case uses one
common JSON instance for all four methods and all 48 performance runs are
optimal. The aggregate statistics are in
\texttt{results/tables/four\_method\_suite\_strengthened\_summary.csv}.

\begin{table}[t]
\centering\scriptsize
\caption{Aggregate statistics over the 12-case A-only performance suite.
Variables and constraints are active scalar Pyomo components; times are in
seconds.}
\label{tab:four_method_suite_summary}
\begin{tabular}{lrrrr}
\toprule
Method & Mean variables & Mean constraints & Mean CPU & Mean wall \\
\midrule
Legacy paper & 14\,225.417 & 26\,923.917 & 1.059189 & 2.414120 \\
Legacy split & 14\,225.417 & 29\,823.917 & 1.478338 & 3.040388 \\
Legacy corrected strengthened & 14\,544.167 & 15\,761.833 & 0.622188 & 1.451844 \\
Event-based optimized & 4\,979.167 & 15\,299.167 & 0.567618 & 1.256824 \\
\bottomrule
\end{tabular}
\end{table}

The expanded suite confirms the focused comparisons: optimized Event-based
has the lowest mean CPU and wall time in this regenerated run, while it uses
approximately 66\% fewer variables than strengthened legacy. Its
constraint count is slightly higher than the legacy-paper count because its
route and timestamp formulations replace compact day-indexed rows; this is a
structural trade-off, not an omitted capacity condition. The C13 loophole
regression has the expected verdict: the paper and split rows accept the fixed
violating pattern, whereas the corrected state rejects it. The clean
regression is optimal under all four methods.

\subsubsection{Strengthened corrected-legacy implementation}

We separately evaluated the equivalent implementation strengthening from
\Cref{subsec:corrected_legacy_strengthening}; it is not substituted into the
four-method table above, so those baseline results remain reproducible. The
comparison uses the same 12 A-only performance cases, three repetitions per
case and implementation, unrestricted CPLEX~12.6, and a 60-second limit. All
36 baseline/strengthened pairs have identical optimal objectives. The clean
regression is also objective-matched, and both implementations reject the C13
loophole instance as infeasible.

\begin{table}[t]
\centering\small
\caption{Corrected-legacy implementation strengthening over 36 paired runs.
Times are means in seconds.}
\label{tab:corrected_strengthening}
\begin{tabular}{lrrrrr}
\toprule
Implementation & Variables & Constraints & Build & CPU & Wall \\
\midrule
Corrected baseline & 14\,544.167 & 25\,261.417 & 0.698785 & 0.659484 & 1.557829 \\
Corrected strengthened & 14\,544.167 & 15\,761.833 & 0.643380 & 0.624975 & 1.465910 \\
\bottomrule
\end{tabular}
\end{table}

The strengthened implementation reduces the mean generated constraint count
by 37.61\%, mean CPU time by 5.23\%, and mean wall time by 5.90\%. It is faster
in 22 of 36 paired CPU measurements and 25 of 36 paired wall measurements.
The variable count is unchanged because both compared implementations contain
the same type-A state variables. These measurements support the
strengthening as an implementation improvement, not as a new mathematical
formulation.

\subsubsection{Optimized Event-based implementation}

We compared the optimized Event-based implementation and the strengthened
corrected legacy implementation on the same 12 performance cases and three
repetitions. All 36 pairs are optimal and objective-matched. The optimized and
baseline Event-based implementations also agree on both regression cases; the
strengthened corrected legacy variant retains its
intended infeasible verdict on the C13 loophole case.

\begin{table}[t]
\centering\small
\caption{Strengthened corrected legacy versus optimized Event-based over 36 paired runs. Times
are means in seconds.}
\label{tab:event_optimization}
\begin{tabular}{lrrrrr}
\toprule
Implementation & Variables & Constraints & Build & CPU & Wall \\
\midrule
Corrected legacy strengthened & 14\,544.167 & 15\,761.833 & 0.640321 & 0.619518 & 1.445313 \\
Event-based optimized & 4\,979.167 & 15\,299.167 & 0.478636 & 0.572310 & 1.261526 \\
\bottomrule
\end{tabular}
\end{table}

Relative to the baseline Event-based implementation, the optimized variant
uses 5.68\% fewer variables and 44.90\% fewer constraints, with 26.79\%
lower mean CPU and 19.58\% lower mean wall time. It is faster in 33 of 36
paired CPU measurements and 34 of 36 paired wall measurements. On this test
grid it also improves on strengthened corrected legacy by 7.62\% in mean CPU
and 12.72\% in mean wall time, while using 65.77\% fewer variables and 2.94\%
fewer constraints. It wins 25 of 36 paired CPU and 25 of 36 paired wall
comparisons against strengthened corrected legacy. These results demonstrate
an implementation-level improvement for this instance family, not universal
runtime dominance.

\subsubsection{Validation scope and required robustness checks}

The A-only results are a controlled formulation comparison, not a final
performance claim. The current evidence uses one deterministic generated case
per $(P,H)$ cell, a single maintenance-capable airport pattern, one set of
maintenance parameters, and one unrestricted CPLEX~12.6 configuration. The
four-method suite also contains two targeted regressions, but these are
correctness checks rather than statistically independent performance samples.
Consequently, the observed Event-based reductions in variables and average
wall time should be interpreted as results for this instance family and solver
configuration only.

A broader validation should repeat the same four-method comparison on multiple
seeds and structurally different instances: hub-and-spoke and mixed-airport
networks, sparse and dense connection graphs, different maintenance capacities,
different A-check thresholds and durations, nonzero valid initial counters, and
instances containing simultaneous departures and tight turn windows. The
solver study should additionally repeat each case under CPLEX with and without
presolve, vary the time limit and warm-start policy, and include at least one
independent MILP solver where the model is supported. Each run should retain
status, optimality gap, objective, active variable and constraint counts, build
time, CPU time, wall time, and peak memory. Objective parity should be checked
before timing comparisons are aggregated. These checks are necessary to
separate formulation effects from instance, solver, and configuration effects.

An earlier cold/warm-start check was performed for the baseline formulations.
Because the published comparison now uses strengthened corrected legacy and
optimized Event-based implementations, those baseline timing values are not
included in the strengthened tables or plots. A refreshed configuration study
with the strengthened implementations remains a robustness follow-up rather
than evidence used for the present performance claims. Attempts with the
available CBC and HiGHS executables, and with the native CPLEX plugin, failed
at the solver interface before producing independent-solver results; they are
retained as diagnostics only.
